\section{Мета}
\label{sec:objective}

Метою цієї лабораторної роботи є дослідження архітектури та криптографічних механізмів IPsec, принципів встановлення 
безпечних мережевих з’єднань, а також особливостей розгортання IPsec у різних мережевих сценаріях. Завдання, які 
ставилися у роботі включають в себе аналіз основних протоколів IPsec (IKEv2, ESP, AH), розробку рішень для безпечного 
підключення для п’яти практичних мережевих сценаріїв, детальне вивчення одного з них та порівняльний аналіз 
IPsec/IKEv2 і WireGuard.

\section{Архітектура та механізми IPsec}
\label{sec:task1}

\subsection*{Призначення IPsec та його місце в моделі мережі}

IPsec (Internet Protocol Security) -- це набір протоколів, що забезпечує криптографічну безпеку -- конфіденційність, 
цілісність, автентифікацію та захист від повторного відтворення — для трафіку на IP-layer traffic~\cite{rfc4301}. 
Він працює мережевому (третьому) рівні моделі OSI, що означає, що він прозоро захищає трафік для всіх загроз, 
розташованих вище над ним.

Наприклад, коли IPsec налаштовано між двома шлюзами, кожна програма -- електронна пошта, веб, запити до баз даних 
будуть захищатися автоматично, без будь-яких змін у самих програмах. Це принципово відрізняється від протоколу 
TLS з першої лаби, який працює на транспортному/прикладнішому рівні та захищає лише конкретну програму, яка його 
використовує.~\cite{cloudflare_ipsec}.

\subsection*{Функції IKEv2 та ESP і їхня взаємодія}

Сам протокол IPsec базується на двох протоколах, які виконують різні ролі~\cite{rfc7296, rfc4303}:

\textbf{IKEv2} (Internet Key Exchange version~2) це \textit{протокол площини керування} (control-plane protocol). 
Він працює через UDP-порти 500/4500 і відповідає за:
\begin{itemize}
    \item Взаємну автентифікацію двох однорангових вузлів (peers) за допомогою сертифікатів, pre-shared keys або 
        методів EAP;
    \item Узгодження криптографічних алгоритмів (шифрування, цілісність, Diffie-Helman group);
    \item Виконання обміну ключами за алгоритмом Діффі–Хеллмана для отримання спільного секретного ключа;
    \item Створення, оновлення ключів та видалення асоціацій безпеки (Security Associations, SAs).
\end{itemize}

\textbf{ESP} (Encapsulating Security Payload) — це \textit{протокол рівня даних} (data-plane protocol), IP protocol 
number~50. Він приймає фактичний трафік користувача та:
\begin{itemize}
    \item Шифрує корисне навантаження за допомогою ключів, узгоджених за допомогою IKEv2 (наприклад, AES-256-GCM);
    \item Аутентифікує кожен пакет за допомогою значення перевірки цілісності (integrity check value, ICV);
    \item Забезпечує захист від повторного відтворення за допомогою вікна порядкових номерів.
\end{itemize}

Їхня взаємодія вібувається наступним чином. Спочатку запускається IKEv2 для встановлення безпечного каналу (IKE~SA), 
а потім відбувається узгодження одного або декількох дочірніх SA. Кожен дочірній SA визначає параметри ESP 
(алгоритм, ключі, селектори трафіку) для фактичного потоку даних. Коли термін дії ключів закінчується, IKEv2 
виконує обмін ключами для генерації нових ключів без переривання трафіку. IKEv2 а-ля як "переговорник"{}, який 
налаштовує захищений канал, а ESP -- "виконавець"{}, який фактично шифрує кожен пакет, що проходить через 
нього~\cite{rfc6071, nist_sp800_77r1}.

\subsection*{Асоціація безпеки (Security Association, SA)}

Асоціація безпеки -- це одностороннє логічне з’єднання, яке містить усе необхідне для обробки захищеного трафіку 
в одному напрямку~\cite{rfc4301}. SA характеризується:
\begin{itemize}
    \item \textbf{SPI} (Security Parameters Index) -- 32-бітний ідентифікатор, що міститься в кожному заголовку 
        ESP;
    \item \textbf{Destination IP address} -- разом із SPI однозначно ідентифікує SA;
    \item \textbf{Protocol} -- ESP (або AH);
    \item \textbf{Cryptographic parameters} -- алгоритм, ключ, довжина вектора ініціалізації (IV), термін дії;
    \item \textbf{Mode} -- транспортний або тунельний;
    \item \textbf{Anti-replay window} стан.
\end{itemize}

Оскільки SA є однонаправленим, то для двостороннього обміну даними потрібно \textbf{щонайменше два SA} -- по одному 
для кожного напрямку. На практиці IKEv2 створює пару дочірніх SA (один вхідний, один вихідний) для кожного захищеного 
потоку даних. Сам IKE~SA (використовується для керуючих повідомлень IKEv2) є окремим двонаправленим SA, тому типова 
сесія IPsec використовує щонайменше три SA загалом~\cite{rfc7296}.

\subsection*{Бази даних SPD та SAD.}

IPsec використовує дві бази даних на кожному хості/шлюзі~\cite{rfc4301}:

\textbf{SPD} (Security Policy Database) -- це \textit{policy table}. Для кожного пакета спочатку йде звернення 
до SPD, і таблиця надсилає відповідь: "Що має статися з цим трафіком?"{}. Далі кожен запис зіставляється за 
селекторами (IP-адреса джерела/призначення, порт, протокол) і визначається одна з трьох дій: \textsc{Protect} 
(застосувати IPsec), \textsc{Bypass} (надіслати у вигляді відкритого тексту) або \textsc{Discard} (відкинути). 

\textbf{SAD} (Security Association Database) -- це \textit{state table}. У ній зберігаються активні SA, тобто 
фактичні ключі, алгоритми, SPI, лічильники послідовностей, які ESP використовує для шифрування та дешифрування 
пакетів. Коли SPD вказує "захистити"{}(protect), система шукає відповідну SA в SAD. Якщо SA не існує, система 
запускає IKEv2 для її створення. Коротше кажучи: SPD визначає, \textbf{що} захищати; SAD містить інформацію про 
те, \textbf{як} це робити~\cite{nist_sp800_77r1}.

\subsection*{Transport vs.\ tunnel mode}

У \textbf{Transport mode} шифрується лише корисне навантаження (payload) IP (заголовок TCP/UDP + дані), а 
оригінальний заголовок IP залишається незмінним~\cite{rfc4303}. Структура пакета стає такою:

\noindent\texttt{[Original IP Header][ESP Header][Encrypted Payload][ESP Trailer][ICV]}

Типове використання цього, це наскрізний захист між двома хостами, що безпосередньо взаємодіють, наприклад, двома
серверами в одній мережі, які обмінюються конфіденційними записами бази даних.

У \textbf{Tunnel mode} шифрується весь оригінальний IP-пакет (заголовок + корисне навантаження) і обгортається 
новим IP-заголовком~\cite{rfc4303}:

\noindent\texttt\scriptsize{[New IP Header][ESP Header][Encrypted: Original IP + Payload][ESP Trailer][ICV]}

Типове застосування: site-to-site VPNs між шлюзами. Внутрішні IP-адреси (приватні LAN адреси) приховані всередині 
зашифрованого тунелю, а зовнішній заголовок містить "gateway-to-gateway"{} адреси. Зовнішній спостерігач бачить 
лише IP-адреси шлюзів і не може визначити, які внутрішні хости взаємодіють між собою.

\subsection*{AH vs.\ ESP і чому використання ESP домінує}

\textbf{AH} (Authentication Header, IP protocol~51) заголовок автентифікації забезпечує цілісність та автентифікацію 
всього IP-пакета, включаючи частини \textbf{зовнішнього} IP-заголовка, але \textbf{нічого не шифрує}~\cite{rfc4302}.

\textbf{ESP} забезпечує як шифрування, так і цілісність/автентифікацію корисного навантаження. Завдяки сучасним 
алгоритмам AEAD (AES-GCM, ChaCha20-Poly1305) ESP виконує обидві функції в рамках однієї операції~\cite{rfc4303}.

\noindent ESP значною мірою замінив AH з трьох причин~\cite{nist_sp800_77r1}:

\begin{enumerate}
    \item ESP забезпечує конфіденційність, чого не може зробити AH;
    \item ESP з алгоритмами AEAD забезпечує цілісність корисного навантаження, порівнянну з AH;
    \item AH перевіряє на автентичність змінні поля IP-заголовка (наприклад, TTL), що спричиняє проблеми перевірки 
        NAT--AH цілісності, коли NAT модифікує IP-заголовок, що робить AH несумісним із середовищами NAT.
\end{enumerate}

\subsection*{NAT і NAT Traversal (NAT-T)}

\textbf{NAT} (Network Address Translation) змінює IP-заголовки та номери портів під час проходження пакетів через
пристрій NAT. Це створює дві проблеми для IPsec~\cite{rfc3948}:

\begin{enumerate}
    \item ESP це IP-level protocol (protocol number~50) без номерів портів TCP/UDP. Багато пристроїв NAT не можуть 
        відстежувати потоки ESP, оскільки вони покладаються на номери портів для розрізнення з’єднань;
    \item Навіть якщо NAT може пропускати ESP, модифікація заголовка IP робить перевірки цілісності AH недійсними 
        та порушує обчислення хеш-сум IKE (hash calculations).
\end{enumerate}

\textbf{NAT-T} (NAT Traversal, RFC~3948)~\cite{rfc3948} вирішує цю проблему таким чином:
1. 
2. Коли NAT виявлено, обидва вузли переходять з UDP 500 на UDP 4500.
3. Кожен пакет ESP упаковується в датаграм UDP 4500: [IP][UDP:4500][ESP][Корисне навантаження].
Тепер пристрій NAT може відстежувати потік за допомогою порту UDP, а вихідний пакет ESP надходить
без змін.
Цей механізм є повністю прозорим — IKEv2 автоматично виявляє та вмикає NAT-T без ручного
налаштування [3].

\begin{enumerate}
    \item Під час узгодження IKEv2 обидва вузли надсилають дані для виявлення NAT (hashes of IP:port). 
        Якщо хеші не збігаються, це означає наявність NAT;
    \item Коли NAT виявлено, обидва вузли переходять з UDP 500 на UDP 4500;
    \item Кожен пакет ESP загортається в UDP~4500 датаграму:
        \texttt{[IP][UDP:4500][ESP][Payload]}. Тепер пристрій NAT може відстежувати потік за допомогою порту 
        UDP, а вихідний пакет ESP надходить без змін.
\end{enumerate}

Цей механізм є повністю прозорим -- IKEv2 автоматично виявляє та вмикає NAT-T без додаткового ручного налаштування~\cite{rfc7296}.

\subsection*{Perfect Forward Secrecy (PFS) в IPsec.}

У контексті IPsec, PFS означає, що ключі, які використовуються для захисту певної сесії передачі даних, неможливо 
вивести з ключів, що використовуються в будь-якій іншій сесії, навіть якщо довгостроковий ключ автентифікації 
(приватний ключ сертифіката або PSK) згодом буде скомпрометований~\cite{rfc7296}. Без ідеальної повної/прямої 
секретності (PFS) зловмисник, який отримав довгостроковий ключ, потенційно може вивести всі сесійні ключі з 
записаних IKE exchanges та розшифрувати весь минулий і майбутній трафік.

Коли PFS увімкнений, IKEv2 виконує новий обмін за алгоритмом Діффі–Хеллмана для кожної перегенерації ключа (rekey) 
для дочірнього SA. Навіть якщо довгостроковий ключ буде скомпрометований, значення DH кожного сеансу є короткочасними 
і вже видаленими, тому минулий трафік залишається захищеним. Увімкнення PFS збільшує обчислювальні витрати (один 
додатковий обмін DH на кожну перегенерацію ключа), але ці витрати є незначними при використання сучасного апаратного 
забезпечення і рекомендуються NIST (The National Institute of Standards and Technology)~\cite{nist_sp800_77r1}.

\section{Сценарії деплою IPsec}
\label{sec:task2}

\textbf{Сценарій 1: Site-to-site connection (два офіси).}

\begin{figure}[!htbp]
    \centering
    \begin{tikzpicture}[
            node distance=1.2cm and 2cm,
            host/.style={draw, rounded corners, minimum width=1.4cm,
                    minimum height=0.7cm, font=\scriptsize\sffamily,
                    fill=blue!8},
            gw/.style={draw, thick, minimum width=1.6cm,
                    minimum height=0.8cm, font=\scriptsize\sffamily\bfseries,
                    fill=orange!15},
            cloud/.style={draw, ellipse, minimum width=2.4cm,
                    minimum height=1cm, font=\scriptsize\sffamily,
                    fill=gray!10},
            net/.style={draw, dashed, rounded corners,
                    inner sep=8pt, gray!60},
        ]
        % LAN A
        \node[host] (h1) {PC-A1};
        \node[host, below=0.5cm of h1] (h2) {PC-A2};
        \node[gw, right=1.5cm of $(h1)!0.5!(h2)$] (gw1) {Gateway A};
        \node[net, fit=(h1)(h2),
            label={[font=\scriptsize]above:LAN A (10.1.0.0/16)}] {};

        % Internet
        \node[cloud, right=2cm of gw1] (inet) {Internet};

        % LAN B
        \node[gw, right=2cm of inet] (gw2) {Gateway B};
        \node[host, right=1.5cm of gw2, yshift=0.35cm] (h3) {PC-B1};
        \node[host, below=0.5cm of h3] (h4) {PC-B2};
        \node[net, fit=(h3)(h4),
            label={[font=\scriptsize]above:LAN B (10.2.0.0/16)}] {};

        % Connections
        \draw[thick] (h1.east) -- (gw1.west);
        \draw[thick] (h2.east) -- (gw1.west);
        \draw[thick] (gw1) -- (inet);
        \draw[thick] (inet) -- (gw2);
        \draw[thick] (gw2.east) -- (h3.west);
        \draw[thick] (gw2.east) -- (h4.west);

        % IPsec tunnel
        \draw[very thick, red!70, densely dashed]
        (gw1.north) to[bend left=25]
        node[above, font=\scriptsize\sffamily, text=red!70!black]
        {IPsec Tunnel (ESP)} (gw2.north);
    \end{tikzpicture}
    \caption{Scenario 1: Site-to-site IPsec VPN}
    \label{fig:s2s}
\end{figure}

\textbf{Кінцеві точки IPsec}: Шлюзи~A та B (gateway) (крайні маршрутизатори/брандмауери кожного офісу).

\textbf{Режим (Mode)}: Tunnel mode -- шлюзи інкапсулюють цілі пакети локальної мережі всередині тунелю IPsec, 
приховуючи внутрішні адреси від Інтернету

\textbf{Захищений трафік}: Весь IP-трафік між підмережами 10.1.0.0/16 та 10.2.0.0/16.

\textbf{Обґрунтування}: Tunnel mode між шлюзами є стандартною "site-to-site"{} VPN схемою, оскільки на окремих 
робочих станціях не потрібно встановлювати програмне забезпечення — шлюзи прозоро обробляють шифрування. 
Внутрішні хости не знають, що використовується IPsec~\cite{rfc4301, nist_sp800_77r1}.

\bigskip
\noindent\textbf{Сценарій 2: Віддалений працівник.}

\begin{figure}[!htbp]
    \centering
    \begin{tikzpicture}[
            host/.style={draw, rounded corners, minimum width=1.6cm,
                    minimum height=0.7cm, font=\scriptsize\sffamily,
                    fill=blue!8},
            gw/.style={draw, thick, minimum width=1.6cm,
                    minimum height=0.8cm, font=\scriptsize\sffamily\bfseries,
                    fill=orange!15},
            cloud/.style={draw, ellipse, minimum width=2.4cm,
                    minimum height=1cm, font=\scriptsize\sffamily,
                    fill=gray!10},
            net/.style={draw, dashed, rounded corners,
                    inner sep=8pt, gray!60},
        ]
        % Remote worker
        \node[host, fill=green!10] (laptop) {Laptop};
        \node[font=\scriptsize\sffamily, below=2pt of laptop]
        {IPsec client};

        % NAT
        \node[draw, right=1.2cm of laptop, minimum width=1cm,
            minimum height=0.6cm, font=\scriptsize\sffamily,
            fill=yellow!15] (nat) {NAT};

        % Internet
        \node[cloud, right=1.5cm of nat] (inet) {Internet};

        % Office
        \node[gw, right=1.5cm of inet] (gw) {VPN Gateway};
        \node[host, right=1.2cm of gw, yshift=0.35cm] (srv1) {Server 1};
        \node[host, below=0.5cm of srv1] (srv2) {Server 2};
        \node[net, fit=(srv1)(srv2),
            label={[font=\scriptsize]above:Office LAN}] {};

        % Connections
        \draw[thick] (laptop) -- (nat);
        \draw[thick] (nat) -- (inet);
        \draw[thick] (inet) -- (gw);
        \draw[thick] (gw.east) -- (srv1.west);
        \draw[thick] (gw.east) -- (srv2.west);

        % IPsec tunnel
        \draw[very thick, red!70, densely dashed]
        (laptop.north) to[bend left=20]
        node[above, font=\scriptsize\sffamily, text=red!70!black]
        {IPsec Tunnel (ESP + NAT-T)} (gw.north);
    \end{tikzpicture}
    \caption{Scenario 2: Remote worker VPN}
    \label{fig:remote}
\end{figure}

\textbf{Кінцеві точки IPsec}: ноутбук співробітника (програмне забезпечення клієнта IPsec, наприклад, 
strongSwan built-in OS client) та gateway VPN в офісі.

\textbf{Режим (Mode)}: Tunnel mode. Весь вихідний пакет шифрується, завдяки чому внутрішні адреси приховуються, 
а ноутбук виглядає як частина офісної локальної мережі.

\textbf{Захищений трафік}: Весь трафік з ноутбука, призначений для офісної мережі (або весь трафік, якщо роздільне
тунелювання вимкнено) є захищеним.

\textbf{Обґрунтування}: Співробітник підключається з ненадійної публічної мережі, найбільш ймовірно, що через домашній 
маршрутизатор, що використовує NAT. IKEv2 автоматично виявляє NAT і вмикає NAT-T (інкапсуляція UDP 4500). Ноутбук 
проходить автентифікацію за допомогою сертифіката або облікових даних EAP і отримує віртуальну IP-адресу з 
офісної підмережі~\cite{rfc7296, rfc3948}. Вуаля і все чудово)

\bigskip
\noindent\textbf{Сценарій 3: Два сервери (host-to-host).}

\begin{figure}[!htbp]
    \centering
    \begin{tikzpicture}[
            host/.style={draw, rounded corners, minimum width=1.8cm,
                    minimum height=0.8cm, font=\scriptsize\sffamily,
                    fill=blue!8},
        ]
        \node[host, fill=green!10] (s1) {Server A};
        \node[draw, ellipse, right=4cm of s1, minimum width=2.4cm,
            minimum height=1cm, font=\scriptsize\sffamily,
            fill=gray!10] (net) {Untrusted Network};
        \node[host, fill=green!10, right=4cm of net] (s2) {Server B};

        \draw[thick] (s1) -- (net);
        \draw[thick] (net) -- (s2);

        % IPsec
        \draw[very thick, red!70, densely dashed]
        (s1.north) to[bend left=20]
        node[above, font=\scriptsize\sffamily, text=red!70!black]
        {IPsec Transport (ESP)} (s2.north);
    \end{tikzpicture}
    \caption{Scenario 3: Host-to-host IPsec (transport mode)}
    \label{fig:h2h}
\end{figure}

\textbf{Кінцеві точки IPsec}: Безпоседеньо сервери~A та сервер~B. 

Режим: транспортний режим — обидва сервери
мають маршрутизовані IP-адреси та спілкуються безпосередньо, тому немає потреби в інкапсуляції IP-заголовка.
Транспортний режим створює менше навантаження. Захищений трафік: весь трафік між IP-адресами двох
серверів (або певною підмножиною протоколів/портів, визначеною в SPD). Обґрунтування: оскільки обидві
кінцеві точки є фактичними хостами, що взаємодіють (без шлюзу між ними), транспортний режим є природним і
більш ефективним вибором. Він шифрує корисне навантаження, зберігаючи оригінальний IP-заголовок, що
дозволяє стандартним політикам маршрутизації та QoS працювати нормально

\textbf{Режим (Mode)}: Transport mode -- обидва сервери мають маршрутизовані IP-адреси та спілкуються безпосередньо, 
тому немає потреби в додатковій інкапсуляції IP-заголовка. Плюч Транспортний режим створює менше навантаження.

\textbf{Захищений трафік}: Весь трафік між IP-адресами двох серверів (або певною підмножиною протоколів/портів, 
задекларованою в SPD).

\textbf{Обґрунтування}: Оскільки обидві кінцеві точки є фактичними хостами, що взаємодіють (no gateway між ними), 
транспортний режим є природним і більш ефективним вибором. Він шифрує корисне навантаження, зберігаючи оригінальний 
IP-заголовок, що дозволяє стандартним політикам маршрутизації та QoS працювати нормально~\cite{rfc4301}.

\bigskip
\noindent\textbf{Сценарій 4: Hub-and-spoke (HQ + multiple branches).} Так звана "зірка"{} -- (головний офіс і 
кілька філій)

\begin{figure}[!htbp]
    \centering
    \begin{tikzpicture}[
            gw/.style={draw, thick, minimum width=1.5cm,
                    minimum height=0.7cm, font=\scriptsize\sffamily\bfseries,
                    fill=orange!15},
            cloud/.style={draw, ellipse, minimum width=2.6cm,
                    minimum height=1.2cm, font=\scriptsize\sffamily,
                    fill=gray!10},
            net/.style={draw, dashed, rounded corners, inner sep=5pt,
                    gray!60, font=\scriptsize\sffamily},
        ]
        % Internet
        \node[cloud] (inet) {Internet};

        % HQ
        \node[gw, above=1.8cm of inet] (hq) {HQ Gateway};
        \node[net, above=0.3cm of hq, minimum width=2.5cm]
        (hqlan) {HQ LAN};

        % Branches
        \node[gw, below left=1.8cm and 2.5cm of inet] (br1)
        {Branch 1 GW};
        \node[net, below=0.3cm of br1, minimum width=2cm]
        (b1lan) {Branch 1 LAN};

        \node[gw, below=1.8cm of inet] (br2) {Branch 2 GW};
        \node[net, below=0.3cm of br2, minimum width=2cm]
        (b2lan) {Branch 2 LAN};

        \node[gw, below right=1.8cm and 2.5cm of inet] (br3)
        {Branch 3 GW};
        \node[net, below=0.3cm of br3, minimum width=2cm]
        (b3lan) {Branch 3 LAN};

        % Lines to internet
        \draw[thick] (hq) -- (inet);
        \draw[thick] (br1) -- (inet);
        \draw[thick] (br2) -- (inet);
        \draw[thick] (br3) -- (inet);

        % IPsec tunnels (dashed red lines)
        \draw[very thick, red!70, densely dashed]
        (hq.south west) to[bend right=10]
        node[left, font=\scriptsize\sffamily,
            text=red!70!black, pos=0.3] {ESP}
        (br1.north);

        \draw[very thick, red!70, densely dashed]
        (hq.south) to
        node[right, font=\scriptsize\sffamily,
            text=red!70!black, pos=0.3] {ESP}
        (br2.north);

        \draw[very thick, red!70, densely dashed]
        (hq.south east) to[bend left=10]
        node[right, font=\scriptsize\sffamily,
            text=red!70!black, pos=0.3] {ESP}
        (br3.north);
    \end{tikzpicture}
    \caption{Scenario 4: Hub-and-spoke IPsec topology}
    \label{fig:hub_spoke}
\end{figure}

\textbf{Кінцеві точки IPsec}: Шлюз штаб-квартири (HQ gateway) виконує роль хаба, в той час як кожен шлюз філії 
встановлює окремий тунель до шлюзу штаб-квартири.

\textbf{Режим (Mode)}: Tunnel mode на кожному каналі -- кожен шлюз філії інкапсулює свій трафік локальної мережі, 
призначений для штаб-квартири (або іншої філії, маршрутизований через штаб-квартиру).

\textbf{Захищений трафік}: Трафік між підмережею кожного філіалу та підмережею штаб-квартири. Трафік між 
філіалами (Inter-branch) проходить через концентратор у штаб-квартирі (HQ hub)

\textbf{Обґрунтування}: Дана модель добре масштабується для організацій із великою кількістю філіалів -- кожному 
філіалу потрібна лише одна конфігурація тунелю (до штаб-квартири), а шлюз штаб-квартири обробляє маршрутизацію 
між філіалами. Це є стандартною архітектурою корпоративної VPN~\cite{nist_sp800_77r1}.

\bigskip
\noindent\textbf{Сценарій 5: Локальна мережа (LAN) і хмарна інфраструктура.}

\begin{figure}[!htbp]
    \centering
    \begin{tikzpicture}[
            host/.style={draw, rounded corners, minimum width=1.4cm,
                    minimum height=0.7cm, font=\scriptsize\sffamily,
                    fill=blue!8},
            gw/.style={draw, thick, minimum width=1.5cm,
                    minimum height=0.7cm, font=\scriptsize\sffamily\bfseries,
                    fill=orange!15},
            cloud/.style={draw, ellipse, minimum width=2.4cm,
                    minimum height=1cm, font=\scriptsize\sffamily,
                    fill=gray!10},
            cld/.style={draw, rounded corners, minimum width=1.4cm,
                    minimum height=0.7cm, font=\scriptsize\sffamily,
                    fill=cyan!8},
            net/.style={draw, dashed, rounded corners,
                    inner sep=8pt, gray!60},
        ]
        % On-prem
        \node[host] (h1) {Workstation};
        \node[host, below=0.5cm of h1] (h2) {DB Server};
        \node[gw, right=1.5cm of $(h1)!0.5!(h2)$] (gw1) {On-prem GW};
        \node[net, fit=(h1)(h2),
            label={[font=\scriptsize]above:On-premises LAN}] {};

        % Internet
        \node[cloud, right=2cm of gw1] (inet) {Internet};

        % Cloud
        \node[gw, right=2cm of inet, fill=cyan!20] (cgw) {Cloud VPN GW};
        \node[cld, right=1.2cm of cgw, yshift=0.35cm] (vm1) {VM / App};
        \node[cld, below=0.5cm of vm1] (vm2) {Cloud DB};
        \node[net, fit=(vm1)(vm2),
            label={[font=\scriptsize]above:Cloud VPC}] {};

        % Lines
        \draw[thick] (h1.east) -- (gw1.west);
        \draw[thick] (h2.east) -- (gw1.west);
        \draw[thick] (gw1) -- (inet);
        \draw[thick] (inet) -- (cgw);
        \draw[thick] (cgw.east) -- (vm1.west);
        \draw[thick] (cgw.east) -- (vm2.west);

        % Tunnel
        \draw[very thick, red!70, densely dashed]
        (gw1.north) to[bend left=25]
        node[above, font=\scriptsize\sffamily, text=red!70!black]
        {IPsec Tunnel (ESP)} (cgw.north);
    \end{tikzpicture}
    \caption{Scenario 5: On-premises to cloud IPsec VPN}
    \label{fig:cloud}
\end{figure}

\textbf{Кінцеві точки IPsec}: Шлюз локальної мережі (on-premises network gateway) та шлюз VPN хмарного провайдера 
(наприклад, AWS VPN Gateway, Azure VPN Gateway, тощо).

\textbf{Режим (Mode)}: Tunnel mode ідентичний до site-to-site VPN, де хмарна VPC виступає в ролі віддаленої сторони.

\textbf{Захищений трафік}: Весь трафік між локальною підмережею та підмережею хмарної VPC.

\textbf{Обґрунтування}: Основні хмарні провайдери (AWS, Azure, GCP) пропонують керовані шлюзи IPsec VPN, які заміняють 
стандартні тунелі IKEv2/ESP~\cite{aws_vpn_docs, azure_vpn_docs}. Це дозволяє організаціям розширювати свою 
приватну мережу в хмару без винесення даних у загальнодоступний Інтернет. Ця конфігурація по суті є VPN site-to-site 
VPN, де однією зі сторін виступає хмарна VPC.

\section{Детальний аналіз сценарію віддаленої роботи}
\label{sec:task3}

Я для себе обрав другий сценарій (малюнок~\ref{fig:remote}) для поглибленого аналізу -- віддалений співробітник 
підключається з публічної мережі до офісної LAN-мережі через IPsec. Як на мене, дуже актуальний варіант)

\textbf{1. Вибір режиму з'єднання: tunnel або transport.}

\textbf{Tunnel mode} Буде правильним вибором, оскільки ноутбук віддаленого співробітника має публічну (або 
NAT-адресу) IP-адресу, яка не входить до офісної мережі. Щоб отримати доступ до внутрішніх ресурсів (файлових 
серверів, баз даних), ноутбук повинен виглядати так, ніби він знаходиться в офісній локальній мережі. Тунельний 
режим дозволяє це зробити: вихідний пакет (з віртуальною офісною IP-адресою як джерелом) інкапсулюється у 
зовнішній пакет (з реальною IP-адресою ноутбука). VPN-шлюз декапсулює пакет і пересилає його в офісну LAN-мережу 
з віртуальною адресою джерела~\cite{rfc4301}. Транспортний режим тут не підійде, оскільки він зберігатиме оригінальний 
IP-заголовок -- публічна IP-адреса ноутбука буде адресою джерела, яку офісна мережа не розпізнає.

\textbf{2. Асоціації безпеки.}

Між ноутбуком та шлюзом VPN встановлюються наступні асоціації безпеки (SA)~\cite{rfc7296}:

\begin{enumerate}
    \item \textbf{IKE SA} (двонаправлена): Захищає control-plane messages IKEv2. 
        Встановлюється під час \\ \texttt{IKE\_SA\_INIT} і \texttt{IKE\_AUTH}.
    \item \textbf{Дочірня SA---outbound}: Створюється при \texttt{IKE\_AUTH}. 
        Використовується ноутбуком для шифрування трафіку, що надсилається до офісу. Ідентифікується за допомогою 
        SPI, обраного шлюзом VPN.
    \item \textbf{Дочірня SA---inbound}: Відповідна SA у зворотному напрямку, яка використовується шлюзом VPN 
    для шифрування трафіку, що надсилається \textit{до} ноутбука. Ідентифікується за допомогою SPI, обраного ноутбуком.
\end{enumerate}

Якщо працівнику потрібно отримати доступ до декількох підмереж із різними політиками безпеки, тоді необхідні додаткові 
пари дочірніх SA, що можуть бути створені за допомогою обмінів \texttt{CREATE\_CHILD\_SA}.

\textbf{3. Роль IKEv2.}

У цьому сценарії IKEv2 виконує такі кроки~\cite{rfc7296}:

\begin{enumerate}
    \item \texttt{IKE\_SA\_INIT} (2 повідомлення): ноутбук і шлюз VPN обмінюються відкритими значеннями DH та
        одноразовими числами (nonces). Вони узгоджують криптографічні алгоритми для IKE SA. Включено корисні 
        дані для виявлення NAT.
    \item \texttt{IKE\_AUTH} (2 зашифровані повідомлення): Ноутбук проходить автентифікацію (за допомогою 
        сертифіката, EAP, або PSK). Шлюз проходить автентифікацію за допомогою свого сертифіката. Проводиться 
        узгодження селекторів трафіку (які підмережі охоплює тунель). Створюється перша пара дочірніх SA. Ноутбуку 
        присвоюється віртуальна IP-адреса (дані конфігурації - (Configuration Payload).
    \item \texttt{CREATE\_CHILD\_SA} (за потреби): Перегенерація ключів IKE~SA або дочірніх (Child~SAs) після 
    закінчення терміну їх дії. Застосовується PFS шляхом виконання нового обміну DH.
\end{enumerate}

\textbf{4. Функції ESP.}

У цьому випадку ESP виконує~\cite{rfc4303}:
\begin{itemize}
    \item \textbf{Шифрування (Encryption)}: Кожен IP-пакет з ноутбука шифрується за узгодженим алгоритмом 
        (наприклад, AES-256-GCM), що запобігає прочитанню трафіку будь-ким у публічній мережі (Wi-Fi у кав’ярні, 
        провайдер (ISP)).
    \item \textbf{Цілісність та автентифікація}: Кожен пакет ESP містить ICV (Integrity Check Value -- значення 
        перевірки цілісності), обчислене на основі зашифрованого корисного навантаження та заголовка ESP 
        (Encapsulating Security Payload). Будь-яке втручання виявляється, і пакет відкидається.
    \item \textbf{Захист від повторного відтворення (Anti-replay)}: Sliding window порядкових номерів не дозволяє 
        зловмиснику перехоплювати та повторно відтворювати старі пакети.
    \item \textbf{Інкапсуляція} (tunnel mode): Весь оригінальний пакет інкапсулюється (wrapped), приховуючи від 
        публічної мережі внутрішні IP-адреси, номери портів і навіть тип використовуваного протоколу.
\end{itemize}

\textbf{5. Що видно зовнішньому спостерігачеві?}

\begin{table}[!htbp]
    \centering
    \begin{tblr}{
            colspec = {X[2,l] X[1,c]},
            row{1} = {font=\bfseries\small, bg=gray!15},
            hlines, vlines,
            rowsep = 3pt,
        }
        Інформація                            & Доступно? \\
        Публічна IP-адреса ноутбука             & Так      \\
        Публічна IP-адреса шлюзу VPN        & Так      \\
        Використовується IPsec/ESP (протокол 50 або UDP 4500)                           & Так      \\
        Обсяг трафіку (розмір пакету)       & Так      \\
        Часові параметри пакетів                     & Так      \\
        Внутрішня (віртуальна) IP-адреса ноутбука    & Ні       \\
        IP-адреса призначення офісних серверів       & Ні       \\
        Application протокол (HTTP, SSH, etc.) & Ні       \\
        Зміст корисного навантаження                       & Ні       \\
        Номери портів TCP/UDP                   & Ні       \\
    \end{tblr}
    \caption{Видимість інформації для стороннього спостерігача}
    \label{tab:visibility}
\end{table}

Підсумовуючи, спостерігач знає, що \textit{ноутбук взаємодіє} з шлюзом VPN і може оцінити, \textit{який обсяг даних} 
передається, але не може визначити, \textit{до чого здійснюється доступ} або \textit{який саме цей вміст}~\cite{rfc4303}.

\textbf{6. NAT та NAT-T.}

У цьому сценарії \textbf{найблільш ймовірне використання NAT} -- віддалений працівник майже напевно підключається через 
домашній маршрутизатор або мобільний хотспот, який використовує NAT (Network Address Translation).

Тому NAT-T (NAT Traversal) є необхідним. Як я згадував вже в розділі~\ref{sec:task1}, IKEv2 автоматично виявляє NAT 
під час \texttt{IKE\_SA\_INIT} і перемикається на інкапсуляцію UDP~4500~\cite{rfc3948}. Цей процес є видним для 
користувача. Без NAT-T пакети ESP (протокол IP~50) або відкидалися б NAT-пристроєм (який не може призначити 
відображення портів для протоколу, що не є TCP/UDP), або порушувалися б, якщо кілька співробітників за одним і 
тим же NAT намагаються підключитися одночасно (виникали SPI колізії).

\textbf{7. Загрози, які усуваються і не усуваються IPsec.}

\textit{Загрози, які усуваються IPsec:}
\begin{itemize}
    \item \textbf{Прослуховування} в публічних мережах Wi-Fi або на каналах інтернет-провайдера -- шифрування 
        ESP унеможливлює зчитування трафіку зловмисником;
    \item \textbf{Фальсифікація пакетів} -- перевірки цілісності ESP виявляють і відхиляють змінені пакети;
    \item \textbf{Атаки повторного відтворення} -- вікно захисту від повторного відтворення (anti-replay window) 
        відхиляє старі пакети;
    \item \textbf{IP spoofing (підробка)} -- шлюз VPN приймає пакети \textit{лише від автентифікованих} однорангових 
        вузлів IPsec;
    \item \textbf{Несанкціонований доступ до мережі} -- лише автентифіковані користувачі з дійсними обліковими 
        даними можуть встановити з'єднання.
\end{itemize}

\textit{Загрози, які не усуваються IPsec:}
\begin{itemize}
    \item \textbf{Compromised endpoint} -- якщо сам ноутбук заражений шкідливим програмним забезпеченням 
        (скомпроментований), IPsec захищає мережевий канал, але зловмисник уже має доступ до відкритого тексту на 
        кінцевому пристрої;
    \item \textbf{Application-layer атаки} -- SQL-ін'єкція, XSS, фішинг відбуваються вище мережевого рівня і є 
        невидимими для IPsec;
    \item \textbf{Аналіз трафіку} -- спостерігач все ще може бачити розміри пакетів та часові закономірності, 
        що може розкрити моделі використання;
    \item \textbf{Відмова в обслуговуванні} -- IPsec не запобігає об’ємним DDoS-атакам;
    \item \textbf{Крадіжка облікових даних} -- якщо облікові дані користувача для VPN вже викрадено (phishing, 
        keylogger), IPsec не може відрізнити законного користувача від зловмисника, який використовує ці дійсні 
        облікові дані~\cite{nist_sp800_77r1}.
\end{itemize}

\section{Порівняння IPsec/IKEv2 vs.\ WireGuard}
\label{sec:task4}

\begin{table}[!htbp]
    \centering
    \begin{tblr}{
            colspec = {X[1.3,l] X[1.2,l] X[1.2,l]},
            row{1} = {font=\bfseries\small, bg=gray!15},
            hlines, vlines,
            rowsep = 4pt,
        }
        Критерій & IPsec/IKEv2 & WireGuard \\
        %
        Архітектура 
            & Складна, модульна: IKEv2 (управління) + ESP (дані), окремі SPD/SAD бази даних
            & Мінімалістична: єдиний протокольний рівень. близько 5к рядків коду ядра \\
        %
        Встановлення ключів 
            & IKEv2: certificates, PSK, EAP; повне узгодження з життєвим циклом SA 
            & Статичні відкриті ключі обмінюються позаканально (out-of-band); 1-RTT Noise IK рукостискання \\
        %
        Криптографічні алгоритми 
            & На вибір: AES-GCM, AES-CBC, ChaCha20, SHA-256/384/512, DH groups 14--21 
            & Фіксовані: ChaCha20-Poly1305, Curve25519, BLAKE2s, SipHash, HKDF \\
        %
        Вибір алгоритму 
            & Повна гнучкість: однорангові вузли узгоджують набори з-поміж запропонованих 
            & Без узгодження: зміна версії протоколу тягне за собою зміну шифрування \\
        %
        NAT traversal 
            & NAT-T виявлено автоматично; UDP~4500 інкапсуляція~\cite{rfc3948} 
            & Native: завжди працює через UDP; NAT є прозорим \\
        %
        Складність конфігурації 
            & Висока: багато опцій, правила (policies), перетворення (transforms), сертифікати PKI 
            & Низька: одна пара ключів + конфігурація однорангового вузла на кожен вузол; мін. конфігураційний файл \\
        %
        Типові сценарії використання 
            & Міжкорпоративні з'єднання site-to-site, віддалений доступ, cloud connectivity, дотримання нормативних вимог 
            & Однорангові тунелі, особисті VPN, мережеві контейнери, overlay networks \\
    \end{tblr}
    \caption{IPsec/IKEv2 vs WireGuard}
    \label{tab:comparison}
\end{table}

\textbf{Придатність до сценаріїв із розділу~\ref{sec:task2}:}

\begin{table}[!htbp]
    \centering
    \begin{tblr}{
            colspec = {X[0.4,c] X[1.5,l] X[0.8,c] X[2,l]},
            row{1} = {font=\bfseries\small, bg=gray!15},
            hlines, vlines,
            rowsep = 3pt,
        }
        \# & Сценарій         & Найкраще            & Обґрунтування                                                                                               \\
        1  & Site-to-site     & IPsec               & Промисловий стандарт; апаратне прискорення; сумісність із шлюзами будь-яких виробників                      \\
        2  & Remote worker    & IPsec або WireGuard & IPsec: корпоративна автентифікація за EAP/cert, OS-native clients; WireGuard: для невеликих teams простіший \\
        3  & Server-to-server & WireGuard           & Компактний, з низькими накладними витратами; легко розгортається на Linux-серверах без використання PKI     \\
        4  & Hub-and-spoke    & IPsec               & Функції підприємств: DMVPN, інтеграція маршрутизації, vendor support, централізована політика               \\
        5  & LAN to cloud     & IPsec               & Хмарні провайдери (AWS, Azure) пропонують вже вбудовані керовані IPsec VPN шлюзи~\cite{aws_vpn_docs}        \\
    \end{tblr}
    \caption{Рекомендовані технології для кожного випадку}
    \label{tab:recommendation}
\end{table}

Підсумовуючи, \textbf{IPsec} краще підходить для корпоративних середовищ та середовищ із обладнанням різних 
виробників (сценарії 1, 4, 5), де критично важливими є сумісність, дотримання нормативних вимог та інтеграція з 
існуючою мережевою інфраструктурою. В той час як, \textbf{WireGuard} вигідно використовувати у розгортаннях, де 
на першому місці стоїть простота (сценарій 3), коли достатньо легкого та швидкого тунелю між відомими кінцевими 
точками, а його фіксована криптографія усуває ризики, пов’язані з неправильною конфігурацією~\cite{wireguard_paper, wireguard_website}. 
Для сценарію~2 підійде будь-який варіант. Вибір залежатиме від масштабу організації та наявної інфраструктури~\cite{nist_sp800_77r1}.

\section{Висновки}
\label{sec:conclusions}

У цій лабораторній роботі я ознайомився і розглянув архітектуру, механізми та практичне впровадження IPsec. Аналіз 
основних його компонентів показав чітке розділення функцій: IKEv2 відповідає за взаємну автентифікацію, 
узгодження ключів та управління життєвим циклом SA, тоді як ESP забезпечує шифрування кожного пакета, захист цілісності 
та захист від повторного відтворення в площині даних. Модель бази даних SPD/SAD забезпечує гнучку структуру політик, 
а два режими роботи (transport and tunnel) призначені для різних мережевих топологій.

П’ять різних ситуацій сценаріїв розгортання показали, як IPsec адаптується до різноманітних мережевих середовищ. 
Site-to-site tunnels (сценарії 1, 4, 5) використовують тунельний режим між шлюзами для прозорого захисту 
всього трафіку між сайтами. Remote access (сценарій 2) передбачає використання NAT-T та призначення віртуальних 
IP-адрес. Host-to-host communication (сценарій 3) виграє від менших накладних витрат у транспортному режимі. 
Детальний аналіз сценарію роботи віддалених співробітників (фактично, багато хто з нас так працює) показав повний 
життєвий цикл: аутентифікацію та обмін ключами за протоколом IKEv2, інкапсуляцію в тунельному режимі ESP, активацію 
NAT-T, а також \textit{чітку межу між тим, що захищає} IPsec (payload confidentiality, integrity, authentication), 
\textit{і тим, що він не може вирішити} (endpoint compromise, application-layer attacks, traffic analysis).

Порівняння IPsec/IKEv2 з WireGuard продемонструвало фундаментальний компроміс між налаштовуваністю IPsec, 
широкою підтримкою виробників та корпоративними функціями, з одного боку, та простотою, продуктивністю та 
мінімальною поверхнею атаки (minimal attack surface) у WireGuard, з іншого. Вибір між ними залежить від умов 
розгортання, а не від абсолютної технічної переваги.
